\section{Заключение}

\subsection{Результаты работы}

В данной работе исследована задача применения методов 
параллелизма к формальной верификации нейронных сетей. 
Основным результатом является реализация и экспериментальная 
валидация неполной верификации (incomplete verification) 
методом CROWN с использованием Tensor Parallelism~--- 
стратегии распределения вычислений, заимствованной 
из обучения больших языковых моделей.

Достигнуты следующие результаты:

\begin{enumerate}
    \item \textbf{Линейное снижение потребления GPU-памяти.} 
    Экспериментально подтверждено, что TP позволяет снизить 
    пиковое потребление памяти пропорционально числу GPU 
    (коэффициент $\approx 1.985$ при двух GPU). Это позволяет 
    верифицировать модели, не помещающиеся в память одного 
    ускорителя.
    
    \item \textbf{Звуковость bounds.} Для всех протестированных 
    моделей TP-реализация вычисляет bounds, являющиеся 
    корректной овер-аппроксимацией: $lb_{\text{tp}} \leq lb_{\text{ref}}$ 
    и $ub_{\text{tp}} \geq ub_{\text{ref}}$. Для моделей 
    с одной парой шардированных слоёв bounds совпадают 
    с эталоном в пределах машинной точности.
    
    \item \textbf{Автоматическое шардирование.} Реализована 
    функция \texttt{tp\_shard\_bounded\_module}, автоматически 
    преобразующая произвольную полносвязную модель (включая 
    загруженную из ONNX) в TP-вариант без ручного 
    переписывания модели.
    
    \item \textbf{Интеграция с \texttt{auto\_LiRPA}.} 
    Реализованы операторы \texttt{BoundLinearTP\_Col} и 
    \texttt{BoundLinearTP\_Row}, корректно распространяющие 
    границы через шардированные слои с учётом зонирования 
    подграфов и \texttt{AllReduce}-синхронизации.
\end{enumerate}

Выявлено также фундаментальное ограничение: для моделей 
с несколькими шардированными зонами промежуточные bounds 
необходимо вычислять методом IBP (а не CROWN), что приводит 
к потере точности bounds, растущей с глубиной модели. 
Это компромисс между масштабируемостью и точностью, 
присущий Megatron-стилю Tensor Parallelism для задачи 
верификации.

\subsection{Направления дальнейшей работы}

\textbf{Полная верификация (Complete Verification).} 
Интеграция TP с методом Branch-and-Bound ($\beta$-CROWN) 
является наиболее значимым направлением. Это потребует 
решения задач синхронизации решений о ветвлении между GPU, 
координации оптимизации параметров $\beta$ 
и распараллеливания Batch BaB. Реализованная в данной работе 
неполная верификация с TP является необходимым фундаментом 
для этой интеграции.

\textbf{Оптимизация $\alpha$-CROWN с TP.} 
Градиентная оптимизация наклонов $\alpha$ совместима 
с TP, поскольку использует тот же backward pass, что 
и CROWN. Интеграция позволила бы получать более точные 
bounds при распределённых вычислениях.

\textbf{Поддержка свёрточных слоёв.} Текущая реализация 
поддерживает только полносвязные (Linear) слои. 
Расширение на свёрточные сети (CNN) потребует реализации 
TP-операторов для \texttt{BoundConv}, что позволит 
верифицировать модели компьютерного зрения.

\textbf{Масштабирование на большее число GPU.} 
Проведённые эксперименты использовали $P = 2$ GPU. 
Валидация при $P = 4, 8$ GPU позволит оценить 
коммуникационные накладные расходы и определить 
практические пределы масштабируемости подхода.
